\paragraph{\texorpdfstring{$S_4$}{S4} 边界退化中的总秩守恒、单标量分类与 Prym 普适性}
在三 cusp 型边界的主分量属数、等型环面秩以及 \(\rho_3'\)-块的 Prym 识别都已明确之后，还可把这些数据进一步压缩为一条更刚性的结构链。

\begin{theorem}[半稳定边界上的总环面秩守恒律]\label{thm:conclusion-s4-boundary-total-torus-rank-conservation}
沿用推论 \ref{cor:xi-s4-boundary-main-component-h10-decomposition} 与定理 \ref{thm:xi-s4-isotypic-jacobian-toric-rank-by-boundary-type} 的三类边界记号 \(r\in\{1,2,3\}\)。定义
\[
g_{\mathrm{main}}(r):=\dim_{\mathbb C}H^{1,0}(X_{\mathrm{main}}),
\qquad
t_{\mathrm{tot}}(r):=t_{\varepsilon}+t_{\rho_2}+t_{\rho_3}+t_{\rho_3'}.
\]
则三类边界同时满足严格守恒律
\[
\boxed{\;
g_{\mathrm{main}}(r)+t_{\mathrm{tot}}(r)=16
\qquad (r=1,2,3).\;}
\]
\end{theorem}
\begin{proof}
由推论 \ref{cor:xi-s4-boundary-main-component-h10-decomposition}，
\[
r=1:\ H^{1,0}(X_{\mathrm{main}})\cong \varepsilon\oplus \rho_3',
\]
\[
r=2:\ H^{1,0}(X_{\mathrm{main}})\cong 2\varepsilon\oplus \rho_2\oplus 2\rho_3',
\]
\[
r=3:\ H^{1,0}(X_{\mathrm{main}})\cong \varepsilon\oplus \rho_2\oplus \rho_3\oplus 2\rho_3'.
\]
利用不可约表示维数
\[
\dim\varepsilon=1,\qquad
\dim\rho_2=2,\qquad
\dim\rho_3=\dim\rho_3'=3,
\]
得到
\[
g_{\mathrm{main}}(1)=1+3=4,
\]
\[
g_{\mathrm{main}}(2)=2+2+2\cdot 3=10,
\]
\[
g_{\mathrm{main}}(3)=1+2+3+2\cdot 3=12.
\]

另一方面，由定理 \ref{thm:xi-s4-isotypic-jacobian-toric-rank-by-boundary-type}，
\[
(t_{\varepsilon},t_{\rho_2},t_{\rho_3},t_{\rho_3'})
=
\begin{cases}
(1,2,3,6),& r=1,\\
(0,0,3,3),& r=2,\\
(1,0,0,3),& r=3.
\end{cases}
\]
故
\[
t_{\mathrm{tot}}(1)=1+2+3+6=12,
\]
\[
t_{\mathrm{tot}}(2)=0+0+3+3=6,
\]
\[
t_{\mathrm{tot}}(3)=1+0+0+3=4.
\]
逐项相加即得
\[
4+12=10+6=12+4=16.
\]
证毕。
\leanverified{paper\_conclusion\_s4\_boundary\_total\_torus\_rank\_conservation}
\end{proof}

\leanverified{paper\_conclusion\_s4\_boundary\_type\_by\_total\_toric\_rank}
\begin{corollary}[边界类型由总环面秩单独决定]\label{cor:conclusion-s4-boundary-type-by-total-toric-rank}
沿用定理 \ref{thm:conclusion-s4-boundary-total-torus-rank-conservation} 的记号，则
\[
\boxed{\;
r=1 \iff t_{\mathrm{tot}}=12,\qquad
r=2 \iff t_{\mathrm{tot}}=6,\qquad
r=3 \iff t_{\mathrm{tot}}=4.\;}
\]
因此，对这三类边界而言，无需读取整个等型秩向量
\[
(t_{\varepsilon},t_{\rho_2},t_{\rho_3},t_{\rho_3'}),
\]
仅凭单一整数 \(t_{\mathrm{tot}}\) 即可完成内禀判别。
\end{corollary}
\begin{proof}
由定理 \ref{thm:conclusion-s4-boundary-total-torus-rank-conservation} 的证明，三类边界对应的总环面秩分别为 \(12,6,4\)，且彼此互异。故反向判别立即成立。证毕。
\end{proof}

\begin{theorem}[\texorpdfstring{$\rho_3'$}{rho3'}-等型因子是唯一的普适边界载体]\label{thm:conclusion-s4-rho3prime-unique-universal-boundary-carrier}
\leanverified{paper\_conclusion\_s4\_rho3prime\_unique\_universal\_boundary\_carrier}
沿用结论 \ref{con:xi-terminal-zm-s4-jacobian-group-algebra-prym-e3-b3} 与结论 \ref{con:xi-terminal-zm-s4-jacobian-a3p-isog-prym3-cubed} 的记号。
记四个 \(\QQ[S_4]\)-等型因子分别为
\[
A_{\varepsilon},\qquad A_{\rho_2},\qquad A_{\rho_3},\qquad A_{\rho_3'}(=A_{3'}).
\]
在文稿既有记号下，可分别识别为
\[
A_{\varepsilon}\sim_{\QQ}J(H),\qquad
A_{\rho_2}\sim_{\QQ}J(Z)^2,\qquad
A_{\rho_3}\sim_{\QQ}J(C_F)^3,\qquad
A_{\rho_3'}\sim_{\QQ}P^3.
\]
则 \(A_{\rho_3'}\sim_{\QQ}P^3\) 是唯一同时满足下列两条性质的等型因子：
\begin{enumerate}
  \item 在每个边界类型 \(r=1,2,3\) 中，对应不可约表示都出现在 \(H^{1,0}(X_{\mathrm{main}})\) 内；
  \item 在每个边界类型 \(r=1,2,3\) 中，对应环面秩都严格为正。
\end{enumerate}
\end{theorem}
\begin{proof}
先逐个比较四类等型在主分量中的出现情形与环面秩表。

对 \(\varepsilon\)-等型，由推论 \ref{cor:xi-s4-boundary-main-component-h10-decomposition} 可知其在三类边界中都出现；但由定理 \ref{thm:xi-s4-isotypic-jacobian-toric-rank-by-boundary-type}，
\[
t_{\varepsilon}=
\begin{cases}
1,& r=1,\\
0,& r=2,\\
1,& r=3,
\end{cases}
\]
故它不满足“环面秩始终非零”。

对 \(\rho_2\)-等型，由
\[
r=1:\ H^{1,0}(X_{\mathrm{main}})\cong \varepsilon\oplus \rho_3'
\]
可见其在 \(r=1\) 时缺失，故不满足“始终出现”。

对 \(\rho_3\)-等型，由推论 \ref{cor:xi-s4-boundary-main-component-h10-decomposition} 可见其在 \(r=1,2\) 时均缺失；并且由定理 \ref{thm:xi-s4-isotypic-jacobian-toric-rank-by-boundary-type}，
\[
t_{\rho_3}=0
\qquad (r=3),
\]
故更不满足所需两条性质。

对 \(\rho_3'\)-等型，则三类边界中的主分量分解分别给出其重数
\[
1,\ 2,\ 2,
\]
故其始终出现；同时对应环面秩
\[
t_{\rho_3'}=
\begin{cases}
6,& r=1,\\
3,& r=2,\\
3,& r=3
\end{cases}
\]
始终严格为正。故唯一性成立。

最后，由结论 \ref{con:xi-terminal-zm-s4-jacobian-a3p-isog-prym3-cubed}，
\[
A_{\rho_3'}=A_{3'}\sim_{\QQ}P^3.
\]
因此所述唯一普适边界载体正是 \(P^3\) 所对应的 \(\rho_3'\)-等型因子。证毕。
\end{proof}

\begin{theorem}[普适边界载体的非主极化刚性]\label{thm:conclusion-s4-universal-boundary-carrier-nonprincipal-polarization}
\leanverified{paper\_conclusion\_s4\_universal\_boundary\_carrier\_nonprincipal\_polarization}
沿用定理 \ref{thm:conclusion-s4-rho3prime-unique-universal-boundary-carrier} 的记号。设
\[
P=\Prym(Y/Z),
\]
并记 \(\lambda_P\) 为其诱导极化。则：
\begin{enumerate}
  \item 由命题 \ref{prop:killo-two-prym-polarization-types}，
  \[
  \mathrm{type}(\lambda_P)=(1,1,1,1,2,2,2,2,2),
  \]
  因而 \((P,\lambda_P)\) 不是主极化阿贝尔簇；
  \item 在结论 \ref{con:xi-terminal-zm-s4-jacobian-a3p-isog-prym3-cubed} 的同源识别
  \[
  A_{\rho_3'}\sim_{\QQ}P^3
  \]
  与自然乘积极化 \(\lambda_P^{\boxtimes 3}\) 的口径下，\(P^3\) 的极化型为
  \[
  (1,1,1,1,2,2,2,2,2)^{\times 3},
  \]
  从而同样不是主极化；
  \item 因此，唯一的普适边界载体 \(A_{\rho_3'}\) 在文稿给定的几何来源下始终携带非主极化。
\end{enumerate}
\end{theorem}
\begin{proof}
主极化的判别条件是极化型全部等于 \(1\)。而命题 \ref{prop:killo-two-prym-polarization-types} 已给出
\[
\mathrm{type}(\lambda_P)=(1,1,1,1,2,2,2,2,2),
\]
其中含有五个 \(2\)，故 \((P,\lambda_P)\) 不是主极化阿贝尔簇。

对乘积极化，极化型按各因子并列排列，因此 \(\lambda_P^{\boxtimes 3}\) 在 \(P^3\) 上诱导的极化型正是
\[
(1,1,1,1,2,2,2,2,2)^{\times 3}.
\]
该型仍含有若干个 \(2\)，故对应乘积极化亦非主极化。

最后，由定理 \ref{thm:conclusion-s4-rho3prime-unique-universal-boundary-carrier}，
唯一的普适边界载体正是与 \(P^3\) 同源的 \(\rho_3'\)-等型因子 \(A_{\rho_3'}\)。因此在文稿既定的 Prym 几何来源下，该普适载体必然携带非主极化。证毕。
\end{proof}

\endinput
